\documentclass[11pt]{article}

\usepackage[T1]{fontenc}
\usepackage{lmodern}
\usepackage{microtype}
\usepackage{amsmath,amssymb,amsthm}
\usepackage[margin=1in]{geometry}
\usepackage[hidelinks,pdftitle={A Uniform Upper Bound for the Diagonal Maximal-Planar Ramsey Number of K1,3 disjoint union K2},pdfauthor={Xiaojian Zeng}]{hyperref}
\usepackage[nameinlink,capitalise]{cleveref}

\newtheorem{theorem}{Theorem}
\newtheorem{lemma}[theorem]{Lemma}
\newtheorem{corollary}[theorem]{Corollary}
\newtheorem{remark}[theorem]{Remark}

\newcommand{\MP}{\mathrm{MP}}
\newcommand{\cK}{\mathcal{K}}
\newcommand{\sqcupdot}{\mathbin{\dot\cup}}

\title{A Uniform Upper Bound for the Diagonal Maximal-Planar\\
Ramsey Number of $K_{1,3}\sqcupdot K_2$}
\author{Xiaojian Zeng\\
\small PBC School of Finance, Tsinghua University, Beijing 100190, China\\
\small \texttt{xiaojianzeng@foxmail.com}}
\date{}

\begin{document}
\maketitle

\begin{abstract}
Let $H=K_{1,3}\sqcupdot K_2$, with the two components on disjoint
vertex sets.  We prove that every red--blue edge-colouring of every finite
simple maximal planar graph on at least $21$ vertices contains a
monochromatic (not necessarily induced) copy of $H$.  Thus the diagonal
maximal-planar Ramsey threshold exists and satisfies
$r_{\MP}(H,H)\leq 21$.  The proof reduces each $H$-free colour graph to
maximum degree at most two or to a four-vertex edge cover, and then uses a
planar bipartite count.  The least threshold is not determined, so the
original exact-value problem remains unresolved.
\end{abstract}

\section{Context and statement}

Several inequivalent Ramsey notions arise when the host graphs are
restricted.  In the weak planar-witness formulation, one asks for a single
planar graph whose every two-edge-colouring contains the target
\cite{AxenovichEtAl2019}; iterated triangulations have also been studied in
that formulation \cite{MaTangYu2020}.  Small Ramsey degree is a different
embedding-colouring invariant: for example, the ordered edge has small
Ramsey degree three in the class of finite linearly ordered split graphs
\cite{Zeng2026SmallRamsey}.  Here we instead use the strong convention of
Baldacchino, Caro and Mifsud \cite{BaldacchinoCaroMifsud2026}: a pair is
unavoidable when every sufficiently large graph in the host class has the
required Ramsey property.

Let $\cK_{\MP}$ be the class of finite simple maximal planar graphs.  For
graphs $A,B$, write $r_{\MP}(A,B)$ for the least integer $N$, when it
exists, such that every red--blue edge-colouring of every
$G\in\cK_{\MP}$ with $|V(G)|\geq N$ has a red copy of $A$ or a blue copy
of $B$.  Copies throughout are ordinary subgraphs, not necessarily induced.
In particular, a copy of
\[
H=K_{1,3}\sqcupdot K_2
\]
consists of a claw and a vertex-disjoint edge of the same colour; extra host
edges among the six selected vertices do not matter.

Baldacchino, Caro and Mifsud classified the connected graphs in the
diagonal maximal-planar problem and asked for a classification of the
disconnected graphs \cite[Theorem~2 and Problem~2]{BaldacchinoCaroMifsud2026}.
The present graph is outside their connected classification.  Our contribution
is the following uniform bound, not an exact evaluation.

\begin{theorem}\label{thm:main}
Let $H=K_{1,3}\sqcupdot K_2$.  Every red--blue edge-colouring of every
$G\in\cK_{\MP}$ with $|V(G)|\geq21$ contains a monochromatic copy of $H$.
Consequently, $r_{\MP}(H,H)$ exists and $r_{\MP}(H,H)\leq21$.
\end{theorem}

The proof is self-contained apart from the standard Euler consequences that
a simple maximal planar graph of order $n\geq3$ has $3n-6$ edges and
minimum degree at least three, and that a simple planar bipartite graph on
$q\geq3$ vertices has at most $2q-4$ edges.

\section{A one-colour reduction}

\begin{lemma}[Claw-cover dichotomy]\label{lem:dichotomy}
If a finite simple graph $F$ contains no copy of $H$, then either
$\Delta(F)\leq2$, or the four vertices of every chosen claw in $F$ form a
vertex cover of $F$.
\end{lemma}

\begin{proof}
If $F$ has no claw, then $\Delta(F)\leq2$.  Otherwise, fix any claw $C$.
An edge with both endpoints outside $V(C)$, together with $C$, is a copy of
$K_{1,3}\sqcupdot K_2$.  Hence every edge meets $V(C)$.
\end{proof}

The dichotomy can be sharpened around a maximum-degree vertex.

\begin{lemma}\label{lem:maxdegree}
Let $F$ be $H$-free and let $v$ have maximum degree $d\geq3$ in $F$.
Then the following hold.
\begin{enumerate}
  \item If $d=3$, the four-set $\{v\}\cup N_F(v)$ covers every edge of
  $F$.
  \item If $d=4$, every edge not incident with $v$ has both endpoints in
  $N_F(v)$.  In particular, every nonisolated vertex of $F$ belongs to
  $\{v\}\cup N_F(v)$.
  \item If $d\geq5$, every edge of $F$ is incident with $v$.
\end{enumerate}
\end{lemma}

\begin{proof}
The first assertion is \cref{lem:dichotomy}, applied to the claw formed by
$v$ and its three neighbours.  Suppose next that $d=4$ and that an edge
$xy$ is not incident with $v$.  If at most one of $x,y$ lies in $N_F(v)$,
then three neighbours of $v$ avoid $x$ and $y$; they form with $v$ a claw
disjoint from $xy$, a contradiction.  Thus both endpoints lie in $N_F(v)$.
Finally, if $d\geq5$, the endpoints of any edge not incident with $v$
exclude at most two neighbours of $v$.  Three remaining neighbours form a
claw with $v$ disjoint from that edge, again a contradiction.
\end{proof}

\section{Proof of the bound}

\begin{proof}[Proof of \cref{thm:main}]
Let $G\in\cK_{\MP}$ have order $n$, and suppose that a red--blue
edge-colouring of $G$ contains no monochromatic $H$.  Let $F_R$ and $F_B$
be the spanning red and blue colour graphs.  Call a colour type L if its
maximum degree is at most two, and type C otherwise.  By
\cref{lem:dichotomy}, every type-C colour contains a claw whose four
vertices cover all edges of that colour.  We bound $n$ in all possible
type pairs.

If both colours are type L, then each has at most $n$ edges.  Hence
\[
3n-6=e(G)=e(F_R)+e(F_B)\leq2n,
\]
which is impossible for $n\geq7$.

Suppose next, after interchanging the colours if necessary, that red is
type L and blue is type C.  Choose a maximum-degree vertex $v$ of $F_B$,
and write $d=d_{F_B}(v)$.

If $d=3$, put $S=\{v\}\cup N_{F_B}(v)$.  By
\cref{lem:maxdegree}, $S$ covers every blue edge.  Every vertex outside
$S$ has degree at least three in $G$ but red degree at most two, so it is
incident with a blue edge.  That edge meets $S$.  Distinct outside
vertices give distinct such cross-edges, while
\[
e_{F_B}(S,V(G)\setminus S)
   \leq \sum_{x\in S}d_{F_B}(x)\leq12.
\]
Therefore $n\leq4+12=16$.

If $d=4$, \cref{lem:maxdegree} says that every vertex outside the
five-set $\{v\}\cup N_{F_B}(v)$ is blue-isolated.  Such a vertex would
have at least three red incident edges, contrary to red being type L.
Thus $n\leq5$.

If $d\geq5$, every blue edge is incident with $v$.  All edges of $G-v$
are consequently red.  Since $d_G(v)\leq n-1$,
\[
e(G-v)=3n-6-d_G(v)\geq2n-5.
\]
On the other hand, the red graph on the $n-1$ vertices of $G-v$ has
maximum degree at most two and hence at most $n-1$ edges.  The inequalities
are incompatible for $n\geq5$.  This completes both mixed type cases.

It remains to consider two type-C colours.  Choose a red claw with
four-vertex set $S_R$ and a blue claw with four-vertex set $S_B$.  Each
set covers all edges of its own colour.  Put
\[
S=S_R\cup S_B,\qquad T=V(G)\setminus S,
\qquad s=|S|\leq8,\quad t=|T|.
\]
There is no edge inside $T$: an edge of either colour would miss the
corresponding cover.  Every vertex of $T$ therefore sends all of its at
least three host edges to $S$.  The simple planar bipartite graph consisting
of the $S$--$T$ edges gives
\[
3t\leq e_G(S,T)\leq2(s+t)-4.
\]
Thus $t\leq2s-4$, and
\[
n=s+t\leq3s-4\leq20.
\]

Every $H$-avoiding colouring therefore has order at most $20$.  This proves
the asserted conclusion for every $n\geq21$.
\end{proof}

\section{Scope and limitations}

The argument proves existence and the explicit upper bound $21$, but gives
no avoiding colouring at order $20$ and no lower bound beyond the trivial
requirement $|V(H)|=6$.  It therefore does not determine the least value of
$r_{\MP}(H,H)$.  Nor does it classify other disconnected graphs in the
broader problem of Baldacchino, Caro and Mifsud.  The comparison above is
limited to the cited literature; no global priority claim is made.

AI assistance was used in literature screening, proof organization,
consistency checking, and editing.  The mathematical argument was checked
separately at the definition and quantifier level.  No exhaustive search,
computer algebra, numerical experiment, or proof assistant is used in the
proof of \cref{thm:main}.

\bibliographystyle{plain}
\bibliography{references}

\end{document}
